\begin{papercard}{mathur2019}
\cf{Stages} \stagetag{C} \stagetag{H} (also \stagetag{D}). A large-scale, automated \emph{measurement} of dark patterns: a crawler visited about 11K shopping sites, text clustering organised about 1.3 million page segments for expert coding, and the study produced a taxonomy of 15 types in 7 categories plus a five-dimension taxonomy of \emph{characteristics}. Its transferable parts are the characteristic dimensions and the coding validation. The crawling is out of scope for an app walkthrough.

\cf{Definition and contribution} Dark patterns are design choices that benefit a service by coercing, steering or deceiving users into decisions they might not make if fully informed \pc{mathur2019}{2}{Dark patterns [32, 48] are user interface design choices that benefit an online service}. Contributions: automated discovery techniques, a prevalence dataset, a descriptive taxonomy of characteristics, and a list of third-party enablers \pc{mathur2019}{3}{We contribute a novel descriptive taxonomy that provides precise terminology}. Prior taxonomies (Brignull, Conti \& Sobiesk, Bösch et al., Gray et al.) were based on anecdotes or submissions, not prevalence \pc{mathur2019}{5}{we provide large-scale evidence documenting the presence and prevalence of dark patterns in the wild}.

\cf{Five characteristic dimensions (the taxonomy)} \textbf{Asymmetric}: unequal weights or burdens on explicit choices (a prominent accept button, a hidden opt-out). \textbf{Covert}: the effect of the design is hidden from users. \textbf{Deceptive}: induces false beliefs through misstatements, misleading statements or omissions. \textbf{Hides information}: obscures or delays necessary information. \textbf{Restrictive}: removes choices \pc{mathur2019}{5}{Does the user interface design impose unequal weights or burdens on the available choices}, \pc{mathur2019}{6}{Hides Information: Does the user interface obscure or delay the presentation of necessary information}. Each type is also linked to cognitive biases: anchoring, bandwagon, default, framing, scarcity and sunk cost \pc{mathur2019}{6}{Default Effect [54]: The tendency of individuals to stick with options that are assigned}.

\cf{Corpus building} Alexa Top Sites (361,102 sites); two classifiers were evaluated against a manual ground truth of 500 random sites (57 shopping), with Webshrinker better (accuracy 94\,\% vs.\ 89\,\%, false-negative rate 18\,\% vs.\ 93\,\%). After a language filter, 19,455 English shopping sites remained \pc{mathur2019}{7}{Compared to Webshrinker, Alexa’s classifications performed poorly on our sample of websites}. A logistic-regression product-URL classifier (83\,\% cross-validated accuracy, 714 labelled URLs) guided the crawler, which stopped after 100 URLs or 15 minutes per site. The crawler was evaluated on 100 sites (86 correct) and the checkout crawler on 100 product pages (66 reached checkout) \pc{mathur2019}{8}{For 86 of those 100 websites, our crawler successfully extracted and returned legitimate product pages}, \pc{mathur2019}{9}{In 66 of the 100 pages, our crawler reached the checkout page successfully}.

\cf{Data collection and clustering} An OpenWPM-based checkout crawler saved page source, screenshots at each state change and HAR files. Pages were split into text-bearing \emph{segments} (visible block elements with text), with a Mutation Summary library capturing dynamic content, giving about 13M segments \pc{mathur2019}{9}{We defined segments as visible HTML elements that contained no other block-level elements}. Preprocessing removed duplicates (to 1.3M). Bag-of-words with 10,133 tokens was reduced by PCA (3 components, $>$95\,\% variance) and clustered with HDBSCAN (min cluster size 10; 10,277 clusters). Clustering was explicitly \emph{to organise segments for an analyst}, not to discover latent constructs \pc{mathur2019}{10}{Our use of clustering was not to discover a set of latent constructs in the data}.

\cf{Expert coding and agreement} Pass 1: one researcher scanned clusters and kept 1,768 of 10,277 as potentially relevant. Pass 2: the team built a shared understanding from the dark-pattern and impulse-buying literature; two researchers independently examined a sample of 200 of the 1,768 clusters, and Cohen's $\kappa=0.74$ was computed on the recorded segments. Disagreements were resolved and one researcher coded the rest, after which the team iteratively grouped patterns into types and categories \pc{mathur2019}{11}{To measure agreement between the researchers, we computed Cohen’s kappa between the segments that were recorded}. Deception checks: sites were revisited every four hours for five days to find countdown timers that reset or offers that remained valid \pc{mathur2019}{11}{we used our checkout crawler to ‘monitor’ the websites containing dark patterns of interest once every four hours}.

\cf{Findings} 1,818 instances on 1,254 of about 11K sites (11.1\,\%), a lower bound because only text on product, cart and checkout pages was examined \pc{mathur2019}{11}{In total, we discovered 1,818 instances of dark patterns from 1,254}. More popular sites had more dark patterns (Spearman $\rho=-0.62$). Seven categories (Table 1): Sneaking (sneak into basket, hidden costs, hidden subscription), Urgency (countdown timer, limited-time message), Misdirection (confirmshaming, visual interference, trick questions, pressured selling), Social proof (activity messages, testimonials), Scarcity (low stock, high demand), Obstruction (hard to cancel), Forced action (forced enrollment). Deception verified in 183 sites (234 instances), e.g.\ 157 deceptive countdown timers on 140 sites \pc{mathur2019}{15}{We considered a countdown timer deceptive if (1) the timer reset after timeout with the same offer still valid}. Trick Questions were counted \emph{only} when the opt-out text began with an affirmative statement about the undesirable practice; Gray et al.'s broader reading was not followed \pc{mathur2019}{17}{We note here that we only considered an opt-out choice as a Trick Question dark pattern when it was misleading}. 22 third-party firms supply dark patterns as a service \pc{mathur2019}{23}{We discovered a total of 22 third-party entities, embedded in 1,066 of the 11K shopping websites}.

\cf{Limitations stated} Text-only patterns, a grey area between deliberate and accidental design, simplified bias mappings, crawler failures, and no coverage of the home, search, account-creation or post-purchase pages \pc{mathur2019}{26}{Our research has several limitations. First, we only take into account text-based dark patterns}.

\cf{Use for this thesis} (1) The five characteristic dimensions are a compact, \emph{observable} scheme to code any suspect interface element in a walkthrough, and they let a finding be described without claiming intent (asymmetric, covert, deceptive, hides information, restrictive). (2) The explicit \emph{operational definitions} and boundary decisions (e.g.\ narrowing Trick Questions) are a model for codebook decision rules, and the footnote shows how a reference taxonomy can be deliberately narrowed. (3) The two-coder validation of a 200-cluster sample with Cohen's $\kappa=0.74$ followed by single coding is a published precedent, with the same caveat as above. (4) The deception checks relied on repeated observation over time, which maps to re-visiting an app at several dates. The shopping context, scale and crawling do not transfer to apps with health data; treat the type list as inspiration only.
\end{papercard}
